% The JumpHMM architecture (spec, slide 4), from fig:model_architecture in the JumpHMM paper
% (HMM-w-jumps-paper/arxiv-paper/sections/main_figures.tex). Input once in the preamble, after vnslides.
% Same structure and math labels as the paper; only the colors and fonts change: the bottom tail is gold, the top
% tail link blue, the Poisson clock the color of the field with an ink border, and the arrows and notes body grey.
% Every state is outlined, with no fill: tinted fills turn muddy on the navy field. The thin lines under the bells mix
% toward the field (vnbg), not toward white.
% The figure runs the width of the slide (user, 2026-10-29), so it is flatter than the paper's: the states sit 4.6
% apart instead of 2, the clock sits lower (y = 3.35, not 5) with its text on one line, the emission note moves
% beside the k-th emission, and the trigger's note sits high on its arrow, above that emission.
% The row of states (user, 2026-10-29: the row was poor): circles large enough that N−1 fits with room, labels at
% the size of the notes, bold dots, arrows with visible heads, and no axis line through the circles; the tail
% braces sit right under their states, and the state-index arrow runs between the two braces.
%
% \vnjumphmmarchitecture draws it at the width of the current line.
\newcommand{\vnjumphmmarchitecture}{%
\resizebox{\linewidth}{!}{%
\begin{tikzpicture}[
  state/.style    = {circle, draw=vnslideink, line width=1.4pt, minimum size=1.8cm,
                     inner sep=0pt, font=\LARGE, text=vnslideink},
  tailbot/.style  = {state, fill=vnbg, draw=vnfigaccent},
  tailtop/.style  = {state, fill=vnbg, draw=vnslidelink},
  midstate/.style = {state, fill=vnbg, draw=vnslideink!70!vnbg},
  dots/.style     = {font=\Huge, text=vnslidebody, yshift=-2pt},
  markov/.style   = {<->, >={Stealth[length=7pt, width=6pt]}, line width=0.9pt, vnslidebody,
                     shorten <=3pt, shorten >=3pt},
  trigger/.style  = {->, >=Stealth, thick, dashed, vnslideink},
  jumparc/.style  = {->, >=Stealth, very thick},
  note/.style     = {font=\LARGE, fill=vnbg, inner sep=2pt},
]
  % the chain: tail states, the k-th state, and dots for the states between
  \node[tailbot]  (s1)   at ( 0.0, 0) {$1$};
  \node[tailbot]  (s2)   at ( 4.6, 0) {$2$};
  \node[dots]     (ld)   at ( 9.2, 0) {$\cdots$};
  \node[midstate] (sk)   at (13.8, 0) {$k$};
  \node[dots]     (rd)   at (18.4, 0) {$\cdots$};
  \node[tailtop]  (sNm1) at (23.0, 0) {$N\!-\!1$};
  \node[tailtop]  (sN)   at (27.6, 0) {$N$};
  \draw[markov] (s1)  -- (s2);
  % the dots sit 2pt low so they look centered on the arrows; the arrows stay on the line through the centers
  \draw[markov] (s2)  -- (ld.west |- s2);
  \draw[markov] (ld.east |- sk) -- (sk);
  \draw[markov] (sk)  -- (rd.west |- sk);
  \draw[markov] (rd.east |- sNm1) -- (sNm1);
  \draw[markov] (sNm1)-- (sN);
  % each state's emission, a bell above it
  \draw[vnfigaccent, thick]
    plot[domain=-0.55:0.55, samples=40, smooth, variable=\t]
    ({ 0.0+\t}, {1.3 + 0.85*exp(-\t*\t/0.04)});
  \draw[vnfigaccent!40!vnbg, thin] (-0.55,1.3) -- (0.55,1.3);
  \draw[vnfigaccent, thick]
    plot[domain=-0.55:0.55, samples=40, smooth, variable=\t]
    ({ 4.6+\t}, {1.3 + 0.85*exp(-\t*\t/0.04)});
  \draw[vnfigaccent!40!vnbg, thin] (4.05,1.3) -- (5.15,1.3);
  \draw[vnslidebody, thick]
    plot[domain=-0.55:0.55, samples=40, smooth, variable=\t]
    ({13.8+\t}, {1.3 + 0.85*exp(-\t*\t/0.04)});
  \draw[vnslidebody!40!vnbg, thin] (13.25,1.3) -- (14.35,1.3);
  \draw[vnslidelink, thick]
    plot[domain=-0.55:0.55, samples=40, smooth, variable=\t]
    ({23.0+\t}, {1.3 + 0.85*exp(-\t*\t/0.04)});
  \draw[vnslidelink!40!vnbg, thin] (22.45,1.3) -- (23.55,1.3);
  \draw[vnslidelink, thick]
    plot[domain=-0.55:0.55, samples=40, smooth, variable=\t]
    ({27.6+\t}, {1.3 + 0.85*exp(-\t*\t/0.04)});
  \draw[vnslidelink!40!vnbg, thin] (27.05,1.3) -- (28.15,1.3);
  % the Poisson clock, its trigger from the k-th state, and the forced runs into each tail
  \node[draw=vnslideink, rounded corners=5pt, fill=vnbg, very thick,
        minimum width=2.8cm, minimum height=1.1cm,
        align=center, font=\LARGE] (clock) at (13.8, 3.35)
        {Poisson clock, prob.~$\epsilon$};
  \draw[trigger] (sk.north) -- (clock.south)
    node[note, pos=0.82, right=5pt, anchor=west, text=vnslidebody] {jump triggered};
  % each forced run: straight out of the clock, straight down outside the tail state, and into it from the side, so
  % the arrow crosses no emission bell
  \draw[jumparc, vnfigaccent, rounded corners=10pt]
    (clock.west) --
    node[note, above, text=vnfigaccent, pos=0.5, yshift=1pt]
      {$K\!\sim\!\mathrm{Pois}(\lambda)$ forced steps}
    (-1.6, 3.35) -- (-1.6, 0) -- (s1.west);
  \draw[jumparc, vnslidelink, rounded corners=10pt]
    (clock.east) --
    node[note, above, text=vnslidelink, pos=0.5, yshift=1pt]
      {$K\!\sim\!\mathrm{Pois}(\lambda)$ forced steps}
    (29.2, 3.35) -- (29.2, 0) -- (sN.east);
  % the tails, braced right under their states, and the state index between them
  \draw[decorate,
        decoration={brace, amplitude=6pt, mirror},
        vnfigaccent, thick]
    (-0.8, -1.0) -- (5.4, -1.0)
    node[midway, below=5pt, vnfigaccent, font=\LARGE]
      {$\mathcal{S}_{-}$\ (bottom tail)};
  \draw[decorate,
        decoration={brace, amplitude=6pt, mirror},
        vnslidelink, thick]
    (22.2, -1.0) -- (28.4, -1.0)
    node[midway, below=5pt, vnslidelink, font=\LARGE]
      {$\mathcal{S}_{+}$\ (top tail)};
  \draw[->, >={Stealth[length=7pt, width=6pt]}, vnslidebody, line width=0.9pt]
    (8.0, -1.5) -- (19.6, -1.5)
    node[midway, fill=vnbg, inner sep=4pt, font=\LARGE, text=vnslidebody] {state index};
  % notes: what the k-th state emits, and the ordinary step
  \node[note, align=right, anchor=east, text=vnslidebody]
    at (13.1, 1.75) {state emission\\[-1pt]
      $\mu_k\!+\!\sigma_k\cdot t_5$};
  \node[note, text=vnslidebody, anchor=south west]
    at (14.75, 0.25)
    {prob.~$1\!-\!\epsilon$};
  % the box around what is drawn, so the slide places the figure by what shows
  \pgfresetboundingbox
  \path[use as bounding box] (-1.75,-2.1) rectangle (29.35,4.3);
\end{tikzpicture}%
}%
}
